\begin{itemize}
    \item Mejorar la precisión espacial y temporal en la localización de fuentes neuronales a partir de señales EEG, en comparación con métodos tradicionales.
    \item Obtener una red neuronal profunda capaz de generalizar frente a distintas configuraciones de ruido, geometrías de cabeza y dinámicas de fuente.
    \item Reducir el error de localización (LE) y el error temporal (TE) en escenarios simulados.
    \item Aumentar las tasas de precisión y sensibilidad en la detección de zonas corticales activas.
    \item Validar cualitativamente la coherencia de los resultados estimados con hallazgos clínicos y modalidades complementarias (fMRI, iEEG).
\end{itemize}